% Ch16 WS3 -- the common-ion effect on solubility. Block 3.
\documentclass{apchem-worksheet}

\apcunitword{Chapter}
\apcunit{16}
\apcunittitle{Solubility Equilibria}
\apcdoctitle{Worksheet 3 \textbullet\ The Common-Ion Effect}
\apcsubtitle{Zumdahl \S16.1 \textbullet\ solubility moves; $K_{sp}$ does not}

\begin{document}
\wsheader[$K_{sp}$: \ce{AgCl} $1.6\times10^{-10}$ \textbullet\ \ce{AgBr}
$5.0\times10^{-13}$ \textbullet\ \ce{BaSO4} $1.5\times10^{-9}$
\textbullet\ \ce{PbI2} $1.4\times10^{-8}$ \textbullet\ \ce{PbCl2}
$1.6\times10^{-5}$ \textbullet\ \ce{CaF2} $4.0\times10^{-11}$
\textbullet\ \ce{MgF2} $6.4\times10^{-9}$ \textbullet\ \ce{Ag2CrO4}
$9.0\times10^{-12}$.]

\problem[]{State, for a salt dissolving in a solution that already contains
one of its ions, what happens to each quantity --- and why:}
\begin{parts}
  \item The solubility: \answerblank[6cm]{decreases}
  \item $K_{sp}$: \answerblank[6cm]{unchanged}
  \item The reason: \answerlines[2]{Adding a product shifts the dissolution
        equilibrium to the left by Le Ch\^atelier, so less solid dissolves.
        $K_{sp}$ is an equilibrium constant and changes only with
        temperature.}
\end{parts}

\problem[]{Predict, without calculating, whether each addition decreases the
solubility of \ce{AgCl}. Answer yes or no and say why:}
\begin{parts}
  \item \ce{NaCl}: \answerblank[6.4cm]{yes --- \ce{Cl-} is a common ion}
  \item \ce{AgNO3}: \answerblank[6.4cm]{yes --- \ce{Ag+} is a common ion}
  \item \ce{KNO3}: \answerblank[6.4cm]{no --- neither ion participates}
  \item \ce{NaBr}: \answerlines[2]{Not by the common-ion effect --- neither
        \ce{Na+} nor \ce{Br-} appears in the \ce{AgCl} equilibrium. (A
        separate reaction to form the even less soluble \ce{AgBr} is
        another matter.)}
\end{parts}

\problem[]{Calculate the solubility of \ce{AgCl}:}
\begin{parts}
  \item in pure water: \answerblank[3.2cm]{$1.3\times10^{-5}$~M}
  \item in \SI{0.10}{\Molar} \ce{NaCl}: \workspace[2cm]
        \keyonly{\begin{apcanswer}$[\ce{Cl-}] \approx 0.10$;
        $1.6\times10^{-10} = (s)(0.10)$;
        $s = \mathbf{1.6\times10^{-9}}$~M\end{apcanswer}}
  \item in \SI{0.020}{\Molar} \ce{AgNO3}: \workspace[2cm]
        \keyonly{\begin{apcanswer}$[\ce{Ag+}] \approx 0.020$;
        $1.6\times10^{-10} = (0.020)(s)$;
        $s = \mathbf{8.0\times10^{-9}}$~M\end{apcanswer}}
  \item By what factor did the \SI{0.10}{\Molar} \ce{NaCl} suppress the
        solubility? \answerblank[3.4cm]{about 7900}
\end{parts}

\problem[]{Calculate the solubility of \ce{BaSO4} in \SI{0.010}{\Molar}
\ce{Na2SO4}, and compare with pure water. \workspace[2.6cm]
\keyonly{\begin{apcanswer}$[\ce{SO4^2-}] \approx 0.010$;
$1.5\times10^{-9} = (s)(0.010)$; $s = \mathbf{1.5\times10^{-7}}$~M.
In pure water $s = 3.9\times10^{-5}$~M, so the sulfate suppressed it by a
factor of about 260.\end{apcanswer}}}

\problem[]{Now a salt with a squared ion. Calculate the solubility of
\ce{PbI2}:}
\begin{parts}
  \item in pure water: \workspace[2cm]
        \keyonly{\begin{apcanswer}$4s^3 = 1.4\times10^{-8}$;
        $s^3 = 3.5\times10^{-9}$;
        $s = \mathbf{1.5\times10^{-3}}$~M\end{apcanswer}}
  \item in \SI{0.10}{\Molar} \ce{NaI}: \workspace[2.2cm]
        \keyonly{\begin{apcanswer}$[\ce{I-}] \approx 0.10$;
        $1.4\times10^{-8} = (s)(0.10)^2$;
        $s = \mathbf{1.4\times10^{-6}}$~M\end{apcanswer}}
  \item By roughly what factor did the iodide suppress it?
        \answerblank[3cm]{about 1100}
\end{parts}

\problem[]{The same salt, the other ion. Calculate the solubility of
\ce{CaF2}:}
\begin{parts}
  \item in \SI{0.10}{\Molar} \ce{NaF}: \workspace[2.2cm]
        \keyonly{\begin{apcanswer}$[\ce{F-}] \approx 0.10$;
        $4.0\times10^{-11} = (s)(0.10)^2$;
        $s = \mathbf{4.0\times10^{-9}}$~M\end{apcanswer}}
  \item in \SI{0.10}{\Molar} \ce{Ca(NO3)2}: \workspace[2.6cm]
        \keyonly{\begin{apcanswer}$[\ce{Ca^2+}] \approx 0.10$ and
        $[\ce{F-}] = 2s$:
        $4.0\times10^{-11} = (0.10)(2s)^2 = 0.40s^2$;
        $s^2 = 1.0\times10^{-10}$;
        $s = \mathbf{1.0\times10^{-5}}$~M\end{apcanswer}}
  \item The two solutions contain the same concentration of common ion, yet
        the effects differ enormously. Explain.
        \answerlines[3]{Fluoride is \emph{squared} in the $K_{sp}$
        expression while calcium is not. Pre-loading the squared ion
        therefore suppresses the solubility far more strongly --- by about
        54,000-fold here, against roughly 22-fold for the calcium.}
  \item In part (b), why does a factor of 4 appear even though calcium is
        the added ion? \answerlines[2]{The fluoride still comes entirely
        from the dissolving \ce{CaF2}, so $[\ce{F-}] = 2s$ and squaring it
        gives $4s^2$. Writing $(0.10)(s)^2$ drops that 4 and makes the
        answer wrong by a factor of 2.}
\end{parts}

\problem[]{Calculate the solubility of \ce{Ag2CrO4} in \SI{0.10}{\Molar}
\ce{AgNO3}. Be careful which ion is squared. \workspace[3cm]
\keyonly{\begin{apcanswer}$\ce{Ag2CrO4(s) <=> 2Ag+(aq) + CrO4^2-(aq)}$,
$K_{sp} = [\ce{Ag+}]^2[\ce{CrO4^2-}]$. The silver is pre-loaded:
$[\ce{Ag+}] \approx 0.10$, $[\ce{CrO4^2-}] = s$.
$9.0\times10^{-12} = (0.10)^2(s)$;
$s = \mathbf{9.0\times10^{-10}}$~M. In pure water it is
$1.3\times10^{-4}$~M.\end{apcanswer}}}

\problem[]{A student writes: ``Adding \ce{NaCl} lowers $K_{sp}$ for
\ce{AgCl}, which is why less dissolves.'' Rewrite the sentence correctly and
say what is still true of the ion product.
\answerlines[4]{Adding \ce{NaCl} does not change $K_{sp}$; it changes the
equilibrium \emph{position}. The added \ce{Cl-} shifts the dissolution
equilibrium to the left, so the solubility falls. At the new equilibrium
$[\ce{Ag+}][\ce{Cl-}]$ still equals $1.6\times10^{-10}$ exactly --- the
chloride is much larger and the silver much smaller than in pure water, and
the product is unchanged.}}

\problem[]{Design question. You want to wash a precipitate of \ce{BaSO4}
with as little loss as possible. Would you rinse with pure water or with
dilute \ce{Na2SO4} solution? Justify.
\answerlines[3]{Dilute \ce{Na2SO4}. The sulfate is a common ion, so it
suppresses the dissolution equilibrium and far less \ce{BaSO4} washes away.
Pure water would dissolve the maximum amount the $K_{sp}$ allows.}}

\begin{spiralstrip}{Chapter 16 \textbullet\ blocks 1--2}
  \item Molar solubility of \ce{AgCl} in water:
        \answerblank[3cm]{$1.3\times10^{-5}$~M}
  \item May you rank \ce{AgCl} against \ce{CaF2} by $K_{sp}$?
        \answerblank[1.8cm]{no}
  \item For \ce{M3X}, $K_{sp} = $ \answerblank[1.8cm]{$27s^4$}
  \item $K_{sp}$ changes only with: \answerblank[2.6cm]{temperature}
  \item Insoluble really means: \answerblank[3.4cm]{very small $K_{sp}$}
\end{spiralstrip}

\end{document}
